% Options for packages loaded elsewhere
\PassOptionsToPackage{unicode}{hyperref}
\PassOptionsToPackage{hyphens}{url}
\documentclass[
]{article}
\usepackage{xcolor}
\usepackage{amsmath,amssymb}
\setcounter{secnumdepth}{-\maxdimen} % remove section numbering
\usepackage{iftex}
\ifPDFTeX
  \usepackage[T1]{fontenc}
  \usepackage[utf8]{inputenc}
  \usepackage{textcomp} % provide euro and other symbols
\else % if luatex or xetex
  \usepackage{unicode-math} % this also loads fontspec
  \defaultfontfeatures{Scale=MatchLowercase}
  \defaultfontfeatures[\rmfamily]{Ligatures=TeX,Scale=1}
\fi
\usepackage{lmodern}
\ifPDFTeX\else
  % xetex/luatex font selection
\fi
% Use upquote if available, for straight quotes in verbatim environments
\IfFileExists{upquote.sty}{\usepackage{upquote}}{}
\IfFileExists{microtype.sty}{% use microtype if available
  \usepackage[]{microtype}
  \UseMicrotypeSet[protrusion]{basicmath} % disable protrusion for tt fonts
}{}
\makeatletter
\@ifundefined{KOMAClassName}{% if non-KOMA class
  \IfFileExists{parskip.sty}{%
    \usepackage{parskip}
  }{% else
    \setlength{\parindent}{0pt}
    \setlength{\parskip}{6pt plus 2pt minus 1pt}}
}{% if KOMA class
  \KOMAoptions{parskip=half}}
\makeatother
\usepackage{graphicx}
\makeatletter
\newsavebox\pandoc@box
\newcommand*\pandocbounded[1]{% scales image to fit in text height/width
  \sbox\pandoc@box{#1}%
  \Gscale@div\@tempa{\textheight}{\dimexpr\ht\pandoc@box+\dp\pandoc@box\relax}%
  \Gscale@div\@tempb{\linewidth}{\wd\pandoc@box}%
  \ifdim\@tempb\p@<\@tempa\p@\let\@tempa\@tempb\fi% select the smaller of both
  \ifdim\@tempa\p@<\p@\scalebox{\@tempa}{\usebox\pandoc@box}%
  \else\usebox{\pandoc@box}%
  \fi%
}
% Set default figure placement to htbp
\def\fps@figure{htbp}
\makeatother
\setlength{\emergencystretch}{3em} % prevent overfull lines
\providecommand{\tightlist}{%
  \setlength{\itemsep}{0pt}\setlength{\parskip}{0pt}}
\usepackage{bookmark}
\IfFileExists{xurl.sty}{\usepackage{xurl}}{} % add URL line breaks if available
\urlstyle{same}
% fallback for those not using the hyperref driver hyperxmp:
\makeatletter
\@ifundefined{xmpquote}{\newcommand{\xmpquote}[1]{#1}}{}
\makeatother
\hypersetup{
  hidelinks,
  pdfcreator={LaTeX via pandoc}}

\author{}
\date{}

\begin{document}

\includegraphics[width=1.58551in,height=1.69167in]{media/image1.png}

\begin{quote}
\textbf{UNIVERSIDAD NACIONAL AUTÓNOMA DE MÉXICO}

\textbf{INSTITUTO DE INVESTIGACIONES EN MATEMÁTICAS APLICADAS Y EN
SISTEMAS}

\textbf{ESPECIALIDAD EN ESTADÍSTICA APLICADA}

\textbf{APLICACIÓN DEL MODELO DE XGRADIENT BOOSTING A LA PREDICCIÓN DE
PRECIOS DE SEGUROS}

\textbf{TESINA}

QUE PARA OBTENER EL GRADO DE:

ESPECIALISTA EN ESTADÍSTICA APLICADA

PRESENTA:

\textbf{ADRIÁN SANDOVAL CORDERO}

DIRECTOR DE LA TESINA:

DRA. CLAUDIA IVONNE JUÁREZ GALLEGOS

\includegraphics[width=1.21438in,height=1.72917in]{media/image2.png}

CIUDAD UNIVERSITARIA, CIUDAD DE MÉXICO, MAYO 2026

\textbf{AGRADECIMIENTOS}

A mi Asesora de Tesina que fue paciente conmigo

A los lectores que me apoyaron con este trabajo:

Dra. Claudia Ivonne Juárez Gallegos

Matro. José Florencio Chávez

Mtra. Leticia Gracia-Medrano Valdelamar

Mtra. Patricia Isabel Romero Mares

Dra. Silvia Ruiz Velasco Acosta

A mi esposa Lizeth Roman por toda su paciencia y ser un gran apoyo.

\textbf{RESUMEN}

En la última década, diversas técnicas de \emph{Machine Learning} han
ganado popularidad en el área Actuarial. Su adopción ha ampliado el
alcance de los análisis tradicionales, al permitir abordar problemas
como la predicción del precio de un seguro ofrecido por la competencia,
con el fin de identificar oportunidades de optimización de tarifas y
apoyar la toma de decisiones.

En este trabajo se compara el desempeño de técnicas tradicionales
(Modelo Lineal y Modelo Lineal Generalizado) con la técnica de
\emph{XGradient Boosting} (un modelo de \emph{Machine Learning}) para
predecir el precio de un seguro con cobertura amplia. En particular, se
incluye:
\end{quote}

\begin{itemize}
\item
  Descripción general de los modelos, de las variables y del problema;
\item
  El esquema de evaluación del desempeño predictivo de los modelos;
\item
  Las conclusiones del análisis;
\item
  La bibliografía consultada.
\end{itemize}

\end{document}
